% adm-foliation.tex — Foliation of spacetime into spacelike hypersurfaces
% Inputable TikZ fragment for fig:adm-foliation
% Requires: tikz with calc, arrows.meta (loaded by ehbook.cls)
%
% Schematic spacetime diagram showing two neighbouring spatial slices
% Sigma_t and Sigma_{t+dt}, the unit normal n^mu, the proper-time gap
% alpha dt, and a tangent vector V^mu.  The gap between slices varies
% across the surface to show that alpha is spatially non-uniform.
%
% Corners of each sheet (perspective parallelogram):
%   Front-to-back perspective offset: dx = -0.7, dy = +1.5
%   Lower: FL=(0.5,0) FR=(7.0,0.20) BR=(6.3,1.70) BL=(-0.2,1.50)
%   Upper: FL=(0.5,2.50) FR=(7.0,3.40) BR=(6.3,4.90) BL=(-0.2,4.00)
%   Gap ranges from ~2.5 (left) to ~3.2 (right).
%
\begin{tikzpicture}[
    >={Stealth[length=5pt]},
    font=\footnotesize,
  ]

  % ── Lower sheet (Σ_t) ─────────────────────────────────────────────

  % Fill
  \fill[EHMarginGray, opacity=0.06]
    plot[smooth, tension=0.6] coordinates {
      (0.5, 0.0) (2.5, -0.08) (4.5, 0.06) (7.0, 0.20)}
    -- (6.3, 1.70) -- (-0.2, 1.50) -- cycle;

  % Front edge (closest to viewer)
  \draw[EHMarginGray, semithick]
    plot[smooth, tension=0.6] coordinates {
      (0.5, 0.0) (2.5, -0.08) (4.5, 0.06) (7.0, 0.20)};

  % Back edge
  \draw[EHMarginGray, thin, opacity=0.4]
    (-0.2, 1.50) -- (6.3, 1.70);

  % Side edges
  \draw[EHMarginGray, thin, opacity=0.4]
    (0.5, 0.0) -- (-0.2, 1.50)  (7.0, 0.20) -- (6.3, 1.70);

  % Coordinate gridlines — transverse (depth 1/3, 2/3)
  \draw[EHMarginGray, very thin, opacity=0.15]
    (0.27, 0.50) -- (6.77, 0.70)
    (0.03, 1.00) -- (6.53, 1.20);

  % Coordinate gridlines — longitudinal (width 1/4, 1/2, 3/4)
  \draw[EHMarginGray, very thin, opacity=0.15]
    (2.13, 0.05) -- (1.43, 1.55)
    (3.75, 0.10) -- (3.05, 1.60)
    (5.38, 0.15) -- (4.68, 1.65);

  % Label
  \node[font=\small, text=EHMarginGray, anchor=west]
    at (7.15, 0.20) {$\Sigma_t$};


  % ── Upper sheet (Σ_{t+dt}) ────────────────────────────────────────

  \fill[EHMarginGray, opacity=0.06]
    plot[smooth, tension=0.6] coordinates {
      (0.5, 2.50) (2.5, 2.44) (4.5, 2.96) (7.0, 3.40)}
    -- (6.3, 4.90) -- (-0.2, 4.00) -- cycle;

  \draw[EHMarginGray, semithick]
    plot[smooth, tension=0.6] coordinates {
      (0.5, 2.50) (2.5, 2.44) (4.5, 2.96) (7.0, 3.40)};

  \draw[EHMarginGray, thin, opacity=0.4]
    (-0.2, 4.00) -- (6.3, 4.90);

  \draw[EHMarginGray, thin, opacity=0.4]
    (0.5, 2.50) -- (-0.2, 4.00)  (7.0, 3.40) -- (6.3, 4.90);

  % Gridlines — transverse
  \draw[EHMarginGray, very thin, opacity=0.15]
    (0.27, 3.00) -- (6.77, 3.90)
    (0.03, 3.50) -- (6.53, 4.40);

  % Gridlines — longitudinal
  \draw[EHMarginGray, very thin, opacity=0.15]
    (2.13, 2.73) -- (1.43, 4.23)
    (3.75, 2.95) -- (3.05, 4.45)
    (5.38, 3.18) -- (4.68, 4.68);

  \node[font=\small, text=EHMarginGray, anchor=west]
    at (7.15, 3.40) {$\Sigma_{t+\mathrm{d}t}$};


  % ── Point P on lower sheet ────────────────────────────────────────
  \coordinate (P) at (3.3, 0.60);
  \fill (P) circle (1.5pt);
  \node[font=\scriptsize, anchor=north east, inner sep=2pt]
    at (P) {$P$};


  % ── Unit normal n^μ (EHJetBlue) ───────────────────────────────────
  %
  % Vertical from P to upper sheet.  At x = 3.3, the upper-sheet
  % surface (same depth fraction) sits at approximately y ≈ 3.38.
  \draw[->, EHJetBlue, very thick, shorten >=2pt]
    (P) -- (3.3, 3.38);
  \node[font=\scriptsize, text=EHJetBlue, anchor=east, inner sep=2pt]
    at (3.25, 3.15) {$n^\mu$};


  % ── Proper-time gap α dt (EHAccretionGold) ────────────────────────
  %
  % Measurement bracket to the right of the normal arrow.
  \draw[EHAccretionGold, thin] (3.50, 0.60) -- (3.70, 0.60);
  \draw[EHAccretionGold, thin] (3.50, 3.38) -- (3.70, 3.38);
  \draw[EHAccretionGold, thin, <->] (3.60, 0.60) -- (3.60, 3.38);
  \node[font=\scriptsize, text=EHAccretionGold, anchor=west, inner sep=1pt]
    at (3.75, 1.99) {$\alpha\,\mathrm{d}t$};


  % ── Tangent vector V^μ (EHAccretionGold) ──────────────────────────
  \draw[->, EHAccretionGold, very thick]
    (P) -- ++(1.5, 0.06);
  \node[font=\scriptsize, text=EHAccretionGold, anchor=south west, inner sep=1pt]
    at (4.65, 0.68) {$V^\mu$};


  % ── Right-angle marker at P ───────────────────────────────────────
  %
  % Small square between n^μ (vertical) and V^μ (horizontal).
  \draw[thin] (3.3, 0.72) -- (3.42, 0.72) -- (3.42, 0.60);

\end{tikzpicture}
